\makeabstract
    {
    Understanding Artificial Photosynthesis with Quantum Dots and Cobalt Porphyrin Catalysts
    }
    {
    Tim Churchill\textsuperscript{1}, Yurii Sulima\textsuperscript{1}, Jacob H. Olshansky\textsuperscript{1}
    }
    {
    \textsuperscript{1} Department of Chemistry, Amherst College, Amherst, MA
    }
    {
    This study aims to understand the photophysics behind the relative performance of several cobalt porphyrin catalysts in quantum dot photocatalysis. Co-H, Co-OMe, Co-CF3, and Co-COOH perform differently in carbon reduction photocatalysis, and also affect electron transfer in a variety of ways. The goal was to investigate whether the differences in electron transfer could explain the photocatalytic performance.
    }
    {
    In this study, quenching experiments were performed on Zinc Selenide Quantum Dots with all four aforementioned catalysts. The change in emission, absorption, and excitation lifetime was observed carefully. Then, catalysts were used in a quantum dot photocatalysis system and carbon reduction was measured using a gas chromatograph. These data were examined in parallel.
    }
    {
    The quenching experiments showed Co-COOH to have the largest impact on electron transfer, as measured through lifetime and emission. Similarly, Co-COOH had the best carbon reduction performance. 
    }
    {
    More data is needed to confirm that the electron transfer has a causal effect on the photocatalysis system. However, if it were proven, this could demonstrate that the electron transfer speed is less important that the binding ability of a given catalyst.
    }

    \newpage
    